\section{Approximation conventions and preliminary reductions}
\label{sec:preliminaries}

\begin{definition}
For \(x,b\in\R\), write \(\ELB(x,b)\) if, for every real \(\nu>b\),
there is \(Q\in\N\), \(Q\ge2\), such that
\[
 \forall p\in\Z\ \forall q\in\N,\qquad
 q\ge Q\ \Longrightarrow\ q^{-\nu}\le\abs{x-p/q}.
\]
\end{definition}

\begin{lemma}\label{lem:endpoint}
If \(\ELB(x,b)\), then \(x\) is irrational. Moreover, for
\(b\ge2\),
\[
 2\le\mu(x)\le b.
\]
\end{lemma}
\begin{proof}
If \(x=a/c\), with \(a\in\Z\), \(c\in\N_{>0}\), use
\(p=ka\), \(q=kc\) with \(k\) arbitrarily large. This contradicts
the positive lower bound in the definition.
For any \(\nu>b\), that definition excludes good reduced approximants
with \(q\ge Q\). There are only finitely many smaller denominators
and finitely many relevant numerators for each of them. Thus no
approximation exponent greater than \(b\) belongs to the defining set
of \(\mu(x)\). Dirichlet supplies the opposite inequality.
\end{proof}

It is enough to contradict the following assertion for each fixed
\(\nu>2\kappa\):
\begin{equation}\label{eq:bad-approximations}
 \text{for every }Q,\quad
 \exists p\in\Z,\ q\in\N,\quad
 q\ge Q,\qquad\abs{x-p/q}\le q^{-\nu}.
\end{equation}
Failure of \(\ELB(x,2\kappa)\) implies this assertion for some such
\(\nu\). Since \(x\ne0\), sufficiently large approximants in
\eqref{eq:bad-approximations} have \(p\ne0\).
They also satisfy \(\abs{p/q}\le |x|+1\).
The integers
\[
 w_i=\lceil\log q_i\rceil
\]
can be made successively as large as desired. In particular, one may
require each \(w_i\) to exceed a threshold depending on the previously
chosen weights. This does not require a prescribed growth rate of the
whole sequence of good approximants.

\begin{lemma}[Root-of-unity reduction]\label{lem:torsion}
In proving Theorem~\ref{thm:compatible}, one may assume either that
the powers \(\xi^j\), \(j\ge0\), are distinct or that \(\xi=1\).
The reduction preserves \(x,F,S,d,s\).
\end{lemma}
\begin{proof}
If \(\xi\) is not a root of unity, equality of two distinct powers
would give a positive power equal to \(1\), since \(\xi\ne0\).
If \(\xi^n=1\) for some \(n\ge1\), replace \((\gamma,\xi)\)
by \((n\gamma,1)\). Equation~\eqref{eq:compatibility} remains true
after raising both sides to the \(n\)-th power.
The new slope is nonzero.
\end{proof}

From now through Section~\ref{sec:parameters}, this normalization has
been performed. Its effect is only on fixed constants.
All choices other than the final degree parameter \(H\) are fixed
before any limit \(H\to\infty\) is taken.
For multi-indices \(\alpha\in\N^m\), put
\[
 |\alpha|=\sum_{i=1}^m\alpha_i,\qquad
 w\alpha=\sum_{i=1}^m w_i\alpha_i,\qquad
 \binom{\alpha}{\beta}=\prod_{i=1}^m\binom{\alpha_i}{\beta_i}.
\]
Coordinatewise inequalities between multi-indices have their usual
meaning. Empty products and empty least common multiples equal \(1\).
